\section{Finding 3: The Edit Classifier Does Not Improve the Detector}
\label{sec:classifier}

\subsection{The classifier is accurate}

On the held-out test split (11,991 query images), the edit classifier's two decisions are:

\begin{table}[ht]
\centering
\begin{tabular}{@{}lrrrr@{}}
\toprule
Decision & Precision & Recall & F1 & Dangerous misses \\
\midrule
Detail destroyed? ($\Rightarrow$ distrust ORB)      & 100\%  & 100\%  & 100\%  & 0.0\% \\
Colour changed? ($\Rightarrow$ distrust hsvHash)    & 98.1\% & 97.0\% & 97.5\% & 3.0\% \\
\bottomrule
\end{tabular}
\caption{Edit classifier decisions on the held-out test split.}
\label{tab:classifier}
\end{table}

A \emph{dangerous miss} is calling a pixelated image ``detail intact'', which would trust ORB
exactly where it fails. This never happens: all 1,800 pixelated test images are recognised,
including the $24\times24$ CryptoPunks where ORB itself is weakest. Raw eight-way accuracy is
lower (89.2\%), almost entirely because an exact copy and an untouched image are identical
pixel for pixel and cannot be told apart from the query alone. Both mean ``trust every
signal'', so the confusion costs nothing.

\subsection{But broken signals are silent, not wrong}

Before using the classifier we measured, rather than assumed, which signal each edit breaks.
Each signal was held at a fixed operating point flagging at most 10\% of non-copies, and for
each edit we recorded both how many copies it catches and how often it fires on the
\emph{wrong} original:

\begin{table}[ht]
\centering\small
\begin{tabular}{@{}lccc r@{}}
\toprule
Edit & aHash & pHash & hsvHash & ORB \\
     & caught / wrong & caught / wrong & caught / wrong & caught \\
\midrule
Flip / rotate           & 5.0 / 0.4  & 2.8 / 0.4  & 77.8 / 7.3 & 81.7 \\
Crop / reposition       & 26.1 / 1.8 & 38.7 / 1.7 & 77.8 / 7.4 & 98.3 \\
Background colour       & 60.7 / 7.7 & 82.7 / 7.5 & 16.1 / 3.8 & 98.9 \\
Colour swap / saturate  & 75.2 / 6.1 & 78.2 / 6.0 & 28.3 / 6.7 & 79.8 \\
Pixelated               & 55.3 / 1.9 & 65.5 / 2.7 & 79.3 / 8.2 & 28.5 \\
Text / logo             & 93.4 / 5.0 & 96.7 / 5.7 & 93.1 / 7.8 & 100 \\
\midrule
Untouched non-copies (false alarms) & 8.2 & 8.1 & 9.2 & 13.0 \\
\bottomrule
\end{tabular}
\caption{Per-signal reliability map (\%). ``Caught'' is the share of that edit's copies the
signal votes for; ``wrong'' is how often it votes for a deliberately wrong original.}
\label{tab:reliability}
\end{table}

Where a signal's detection collapses, its false alarms collapse with it. aHash catches only
5.0\% of rotated copies, and it fires on the wrong original only 0.4\% of the time. hsvHash on
background changes catches 16.1\% and misfires 3.8\%, below its normal 9.2\%. A broken signal
here does not vote randomly. It abstains. In a two-vote rule, telling the detector to ignore a
signal that is already silent changes nothing.

\subsection{Measured end to end}

We compared the detector with ORB, with and without the classifier, at the same operating
points as Table~\ref{tab:whole}:
\begin{itemize}
  \item static two-of-four with ORB: F1 87.5 (precision 97.8\%, recall 79.2\%);
  \item classifier-guided, no fallback: F1 87.4;
  \item classifier-guided with the static fallback: F1 87.5, identical to static.
\end{itemize}
The swap accounts for the whole gain; the classifier adds nothing measurable on top of it.

The classifier could only help through the other lever: lowering the quorum. When two signals
are distrusted at once, the vote falls to ``one of the remaining two''. With one edit per
image, the classifier almost never distrusts two signals at once. A simulated third flag shows
the lever is real. Distrusting aHash and pHash on rotated images leaves hsvHash and ORB and
lowers the quorum to one, which raises flip/rotate recall from 58.4\% to 97.7\% at an
unchanged 8.5\% false-alarm rate. But that flag's trigger reads transparent corners left by
our generator's rotation, and it scores 0\% on the authors' images. It cannot be relied on, so
it is not counted as a result.

\subsection{The classifier is specific to our generator}

On the authors' reference set the classifier degrades badly. Their CryptoPunks are already
enlarged RGB images, ours are native $24\times24$ RGBA, so the transparency and histogram
``combing'' cues it relies on disappear. It then over-predicts pixelation, and the colour
decision drops to 0\% recall. Over-predicting pixelation only costs recall, because it makes
the detector ignore ORB more often; it never makes the detector trust ORB wrongly. The
classifier's accuracy above is therefore a result within our own generator, not evidence
that it recognises edits in general.

\subsection{The vote rule depends on how common copies are}

Sweeping the quorum on our test set suggests that ``one of four'' beats the paper's ``two of
four'' (F1 94.3 against 87.5). That is an artifact of a test set that is 82.6\% copies.
Reweighting the same results to a realistic share of copies reverses it
(Figure~\ref{fig:prevalence}): at 2\%, 10\% and 25\% copies, two of four is better. Real
copymints are rare among all mints, so the paper's rule is the right one.

\begin{figure}[ht]
\centering
\includegraphics[width=0.66\linewidth]{\fig{fig_prevalence.pdf}}
\caption{F1 of the ORB detector under the two vote rules, reweighted to an assumed share of
copies. The paper's two-of-four rule wins wherever copies are uncommon.}
\label{fig:prevalence}
\end{figure}

\subsection{Phase E: two edits on one image}

If the classifier fails only because single-edit data never breaks two signals at once, then
images with two edits should change the verdict. We generated every ordered pair of the six
non-exact edits (30 combinations; 18,000 copies and 2,400 non-copies, 88.2\% copies). This
test used the original panel with sHash, which stores nothing extra. The classifier would
distrust sHash on pixelation (where sHash drops from 74.3\% on crops to 38.1\% while staying
silent) and hsvHash on a colour change, lowering the quorum to one of two.

\begin{table}[ht]
\centering
\begin{tabular}{@{}lrrrr@{}}
\toprule
Detector (two-edit test set) & Precision & Recall & F1 & $\Delta$F1 \\
\midrule
Paper's panel (with sHash), static     & 96.7\% & 36.4\% & 52.9\% & --- \\
Same panel, classifier-guided          & 96.7\% & 36.4\% & 52.9\% & +0.01 \\
ORB in place of sHash, static          & 97.3\% & 45.1\% & 61.6\% & +8.7 \\
\midrule
Flag everything                        & 88.2\% & 100\%  & 93.8\% & \\
\bottomrule
\end{tabular}
\caption{Phase E results. Precision is about equal across detectors, so recall is the fair
comparison.}
\label{tab:phasee}
\end{table}

The verdict did not change (Table~\ref{tab:phasee}). The reason is physical: \textbf{pixelation
erases the colour trace.} To test whether the colour decision was merely being crowded out by
pixelation in the eight-way classifier, we trained two independent yes/no detectors, one for
destroyed detail and one for a colour change. The colour detector fires on 76--100\% of
colour edits combined with any non-pixelation edit, and on 0\% of pixelation combined with a
colour edit, in both orders. Pixelation collapses the palette, so the histogram combing the
detector reads is gone. The one situation where the classifier could lower the quorum, detail
destroyed \emph{and} colour changed, cannot be observed with these cues. Without the static
fallback the classifier even hurts (F1 50.3), because it drops sHash's occasional correct vote
without ever reaching the quorum change. Meanwhile the ORB swap still pays, at +8.7 F1.

This conclusion is limited to our generator and our pixelation strength, which is severe
(down to 4--12\% of the original size and back). Milder pixelation might leave a colour trace.
We did not test that.
